\chapter{Tensor Lowering and Recurrence}

Tensor lowering is the main compiler path that makes \einlang{} different from
a plain expression interpreter. The source program exposes named axes and
reductions; the compiler turns those source facts into lowered loops, reduction
kernels, recurrence order metadata, and storage decisions.

\section{Einstein Grouping}

Multiple clauses for the same indexed binding are first grouped into a single
\texttt{EinsteinIR}. This is needed for recurrences and piecewise tensor
definitions:

\begin{lstlisting}[style=einlang]
let y[0] = init0;
let y[1] = init1;
let y[t in 2..T] = y[t - 1] + y[t - 2];
\end{lstlisting}

The grouping pass keeps these clauses as one semantic binding instead of three
unrelated assignments. Later passes can therefore see that the final clause
reads the same binding it writes.

The example is short, but it is the reason grouping must run before most tensor
analyses. Without grouping, the first two clauses would look like independent
scalar assignments and the third would look like a separate indexed tensor
definition. Grouping gives the compiler a single object with boundary coverage,
recurrent domain, and self-read offsets. That one object is what range
analysis, recurrence detection, storage inference, and the backend agree on.

\section{Rest Patterns and Implicit Ranges}

Named rest indices such as \texttt{..batch} express a variable number of leading
or trailing dimensions. The rest-pattern preprocessing pass expands them into
compiler-generated index variables once enough rank information is available.
The implicit range detector recovers ranges from shaped reads, reductions, and
index expressions. Together these passes let source programs write compact
generic tensor code while the backend receives explicit loops.

For example, a function that preserves batch axes can be written against an
unknown-rank prefix. The compiler expands that prefix only after rank
information is known, replacing the rest index with concrete generated index
variables. This avoids requiring a separate family of functions for every batch
rank.

\subsection{Rest Expansion Algorithm}

The rest-pattern pass follows a determination-first rule. A rest pattern may be
used in a mixed access only after some access determines how many dimensions it
spans. This prevents ambiguous programs such as two unknown rest patterns
appearing together before either one has a rank.

\begin{lstlisting}[style=pythoncode,caption={Named rest-index preprocessing.}]
expand_rest_patterns(einstein_binding):
    clauses = einstein_binding.clauses
    require all clauses to have the same output rank

    output_rests = rest names in first clause indices
    if output_rests is empty:
        return

    accesses = rectangular accesses inside first clause body
    for rest_name in output_rests:
        require rest_name appears in at least one body access
        require at least one access determines rest_name alone

    rest_dim_mapping = {}
    for rest_name in output_rests:
        access = first body access using rest_name
        array_rank = rank_of_access_array(access) or len(access.indices)
        explicit_count = number of non-rest indices in access
        span = max(0, array_rank - explicit_count)
        rest_dim_mapping[rest_name] = [0, ..., span - 1]

    for clause in clauses:
        new_output_indices = []
        old_rest_defid_to_new_defids = {}
        for index in clause.indices:
            if index is IndexRestIR("batch"):
                for dim in rest_dim_mapping["batch"]:
                    did = resolver.allocate_for_local()
                    new_output_indices.append(IndexVarIR("batch.dim", did))
                    old_rest_defid_to_new_defids[index.defid].append(did)
            else:
                new_output_indices.append(copy index)

        clause.indices = new_output_indices
        clause.value = replace_body_rest_indices(
            clause.value, rest_dim_mapping, old_rest_defid_to_new_defids)

        for generated index used in an array access:
            clause.variable_ranges[index.defid] = 0 .. array.shape[axis]
\end{lstlisting}

The same expansion machinery is used for reduction rests such as
\texttt{sum[..batch](...)}. The result is that range analysis and lowering do
not need a separate representation for unknown-rank groups; they receive
ordinary generated index variables and shape-derived ranges.

The pass is careful to require a determining access because otherwise a rest
name would be underconstrained. In an expression with two unknown-rank tensors,
the compiler cannot safely infer how many axes \texttt{..batch} spans unless at
least one read ties the rest pattern to a known or inferred rank. Once the span
is known, however, generated indices behave exactly like user-written indices:
they receive DefIds, ranges, and later loop structures. This keeps the rest
syntax ergonomic without creating a second class of index variable.

The replacement map from old rest DefIds to generated DefIds is also important
for diagnostics and lowering. The source program has one rest name, but the
lowered compiler needs one loop variable per concrete axis. Keeping a mapping
between those identities lets the compiler report errors against the source
rest while executing explicit generated loops. That is the pattern used
throughout the implementation: preserve the source fact for messages, but make
the backend-facing representation explicit.

The expansion pass is also where rest syntax stops being magical. Before this
point, \texttt{..batch} is a compact source phrase whose arity depends on rank
information elsewhere in the expression. After this point, every generated
dimension is an ordinary index variable with its own DefId and range. That means
range analysis, shape analysis, lowering, and diagnostics do not need special
cases for variable-arity source syntax. They see the same kind of explicit
index objects they would have seen if the programmer had written every axis by
hand.

The determination rule is conservative because ambiguous rest expansion can
silently change program meaning. If two operands both contain \texttt{..batch}
but neither operand has a known rank, there are several possible spans that
could make the syntax type-check. Choosing one would bake an arbitrary layout
decision into the compiler. Requiring a determining access turns that ambiguity
into a source diagnostic and gives the programmer a clear repair: add a range,
an annotation, or a shaped read that fixes the rank. This keeps generic tensor
code flexible without making rank inference guessy.

\textit{Concrete example.} Consider \texttt{let y[..batch, c] =
sum[i](x[..batch, i] * w[i, c]);}. If shape analysis knows that \texttt{x} has
rank three, then \texttt{..batch} spans the first two axes and is expanded into
two generated indices, for example \texttt{batch.0} and \texttt{batch.1}. Those
indices receive ranges from the first two axes of \texttt{x}; \texttt{i}
receives the final axis of \texttt{x} and the first axis of \texttt{w}; and
\texttt{c} receives the second axis of \texttt{w}. If the rank of \texttt{x}
is unknown, the same source is rejected until another annotation or access
determines the rest span.

\section{Einstein Lowering}

\pathcode{src/einlang/passes/einstein_lowering.py} lowers high-level tensor IR
into explicit lowered nodes:

\begin{itemize}
\item \texttt{LoweredEinsteinIR} represents an indexed tensor binding as one or
more lowered clauses.
\item \texttt{LoweredReductionIR} represents reductions such as \texttt{sum},
\texttt{max}, \texttt{min}, \texttt{prod}, \texttt{all}, and \texttt{any}.
\item \texttt{LoweredSelectAtArgmaxIR} represents selected values at extremum
indices.
\item \texttt{LoweredComprehensionIR} represents array comprehensions after
range and constraint analysis.
\end{itemize}

The lowered form keeps loop structures, guards, value bindings, reduction
ranges, and body expressions explicit. The backend therefore executes a checked
loop plan rather than rediscovering source syntax.

\subsection{Lowered Clause Shape}

A lowered clause contains four important pieces: the write indices, the loop
structures that enumerate those indices, value bindings introduced by
\texttt{where}, and guard predicates that decide whether the clause applies.
The backend can evaluate base clauses, guarded clauses, and recurrent clauses
with one common representation. This is why the grouping and constraint passes
run before lowering: by the time the backend sees a clause, the compiler has
already classified which parts are iteration, binding, and guard logic.

\subsection{Lowering Algorithm}

Lowering turns declarative clauses into loop plans while preserving source
structure. It does not execute anything; it makes execution explicit.

\begin{lstlisting}[style=pythoncode,caption={Lowering an Einstein binding.}]
lower_einstein_binding(binding):
    lowered_items = []
    output_shape = union_shape_from_clause_ranges(binding.clauses)

    for clause in binding.clauses:
        loops = []
        bindings = []
        guards = []

        for output_index in clause.indices:
            if output_index is index variable:
                range = clause.variable_ranges[output_index.defid]
                loops.append(LoopStructure(output_index, range))
            else if output_index is literal:
                remember literal write position

        for where item in classified_where_clause(clause):
            if item is value binding:
                bindings.append(lower_binding(item))
            else if item is predicate:
                guards.append(GuardCondition(item))
            else if item is range domain:
                attach to matching loop variable

        body = lower_expr(clause.value)
        lowered_items.append(
            LoweredEinsteinClauseIR(indices=clause.indices,
                                    loops=loops,
                                    bindings=bindings,
                                    guards=guards,
                                    body=body))

    return LoweredEinsteinIR(items=lowered_items,
                             shape=output_shape,
                             element_type=inferred_element_type)

lower_reduction(reduction):
    loops = [LoopStructure(var, reduction.loop_var_ranges[var.defid])
             for var in reduction.loop_vars]
    body = lower_expr(reduction.body)
    return LoweredReductionIR(operation=reduction.operation,
                              loops=loops,
                              body=body,
                              guards=lowered_where_guards)
\end{lstlisting}

Several later optimizations rely on this representation. A reduction nested
inside a clause remains visibly nested; a guard remains a guard rather than a
Python \texttt{if}; and literal-index boundary clauses remain distinguishable
from loop-bearing recurrence clauses.

The lowering algorithm is the point where declarative source clauses become a
backend plan. Output indices are split into literal write positions and loop
variables, where items are split into bindings and guards, and reductions are
converted into explicit loop-bearing nodes. The body expression is lowered only
after that structure is available, so nested reductions and guarded bodies keep
their meaning. This is why earlier classifier passes matter: lowering should
not guess whether a where-clause item is a loop domain or a predicate.

The returned \texttt{LoweredEinsteinIR} is still not executable code in the
bytecode sense. It is an explicit data structure describing writes, loops,
guards, and bodies. That representation gives the NumPy backend freedom to
choose vectorized or scalar execution without changing semantics. It also gives
later recurrence and execution-fact passes a uniform object to inspect; they do
not need to understand the entire high-level syntax again.

The lowered representation also separates allocation from assignment. A binding
may have several clauses, some with literal output positions and some with loop
variables. Lowering computes the shared output shape once, then records each
clause as a write plan into that allocation. This is why boundary clauses, guard
clauses, and interior recurrent clauses can coexist without being flattened into
unstructured Python control flow. The backend receives a compact execution plan,
but the plan still remembers which source clause produced each write.

The reduction branch follows the same philosophy. A reduction is not evaluated
during lowering, and it is not erased into a generic call. It becomes a
\texttt{LoweredReductionIR} with explicit loop structures, guards, and a lowered
body. That makes specialized kernel recognition possible later: the execution
facts pass can inspect a structured reduction and decide whether it is a matrix
contraction, a scalar accumulation, or a nested reduction requiring a safer
path. The value of the algorithm is therefore not only that it creates loops,
but that it creates loops with analyzable intent.

\textit{Concrete example.} The source binding \texttt{let pos[i] = x[i] where
x[i] > 0;} lowers to one loop over \texttt{i}, one guard
\texttt{x[i] > 0}, and a body \texttt{x[i]}. The guard remains a
\texttt{GuardCondition} rather than becoming a host-language \texttt{if} hidden
inside the backend visitor. For \texttt{let fib[0] = 0; let fib[n in 1..N] =
fib[n - 1] + 1;}, lowering produces a literal write clause for \texttt{0} and
a loop-bearing recurrent clause for \texttt{n}. The backend can therefore
initialize the boundary point once and run the recurrent clause in dependency
order.

\section{Backend Execution Facts}

\pathcode{src/einlang/passes/lowered_execution_facts.py} computes summaries for
lowered reductions and clauses. The facts include nested reduction presence,
conditional bodies, loop DefIds, static loop ranges, call arguments that depend
on loop variables, reduction dimension counts, direct nested lowerings, and
candidate vectorization strategies. This is a deliberate compiler/backend
boundary: the compiler owns analysis, while the backend uses attached facts to
choose execution paths.

\subsection{Execution-Fact Algorithm}

Execution facts summarize lowered IR once, in the compiler, so the backend does
not rediscover dependencies by ad hoc inspection at runtime.

\begin{lstlisting}[style=pythoncode,caption={Lowered execution-fact construction.}]
build_lowered_execution_facts(program):
    collector = SummaryCollector()
    collector.collect(program)
    collector.collect(each function in function_ir_map)

    for reduction in collector.lowered_reductions:
        summary = collector.collect(reduction.body)
        plan = recognize_reduction_kernel(reduction)
        strategy = recognize_reduction_strategy(reduction, summary, plan)
        reduction.execution_facts_id = fresh_id()
        reduction.execution_strategy = strategy
        store ReductionFacts(id, nested_reductions, ifs, defids, strategy)
        if plan exists:
            reduction.kernel_plan_id = plan.id

    for clause in collector.lowered_clauses:
        body = collector.collect(clause.body)
        loop_defids = [loop.variable.defid for loop in clause.loops]
        call_arg_loop_defids = body.call_arg_defids intersect loop_defids
        static_ranges = [static_range(loop.iterable) for loop in clause.loops]
        strategy, scalar_dims = recognize_clause_strategy(clause, body)
        clause.execution_facts_id = fresh_id()
        clause.vectorization_strategy = strategy
        clause.vectorization_scalar_loop_dims = scalar_dims
        store ClauseFacts(id, loop_defids, static_ranges, body_summary, strategy)
\end{lstlisting}

The backend treats a missing strategy as a compiler error. That design makes
analysis omissions visible during testing rather than silently choosing a
different execution path.

Execution facts are the compiler/backend contract for tensor performance. The
collector walks lowered IR once and records the facts that would be expensive
or fragile for the backend to rediscover repeatedly: loop DefIds, static
ranges, nested reductions, call dependencies, and vectorization candidates. The
backend then uses an attached ID or strategy rather than performing ad hoc
structural analysis while executing every clause.

This separation improves correctness as much as speed. If a clause body
contains a function call that depends on a loop variable, the compiler can mark
that fact and force a call-scalar path. If a reduction matches a known kernel,
the compiler can attach a plan. If no plan is recognized, tests can detect the
missing strategy immediately. In other words, the algorithm turns implicit
backend assumptions into explicit IR metadata.

The fact IDs are a small but useful indirection. They let lowered nodes carry a
stable handle while the full summaries live in compiler-owned tables. That keeps
IR nodes from becoming oversized containers and lets tests inspect summaries
directly. It also means a backend can fail loudly if a lowered node has no fact
ID even though the pipeline promised one. The absence of metadata becomes an
implementation error, not a silent request to recompute analysis at runtime.

The body summary is especially important for mixed execution strategies. A
clause can look vectorizable at the loop level but contain a function call whose
arguments depend on a loop DefId. Broadcasting such a call would change the
calling convention or produce an object array. By recording call dependencies,
nested reductions, conditionals, and static ranges, the compiler can choose a
hybrid or scalar path only where it is needed. The backend optimization is then
bounded by facts that were computed once in a phase designed for structural
inspection.

\textit{Concrete example.} A lowered matrix multiplication reduction receives a
summary with one reduction loop, no conditionals, no loop-dependent calls, and a
body that is a product of two rectangular accesses; the recognizer can attach a
specialized reduction strategy. A clause such as \texttt{let y[i] = f(i)}
instead records that a call argument depends on the loop DefId for \texttt{i},
so the backend chooses the call-scalar path. Both programs may be small, but
the attached facts are different enough that the backend should not rediscover
them by inspecting Python objects during execution.

\section{Reduction and Vectorization Paths}

The NumPy backend has several execution strategies for lowered tensor nodes:
fully vectorized clauses, scalar loops, hybrid clauses, call-scalar paths, and
specialized reduction kernels. Matrix-like contractions can be recognized and
sent through efficient NumPy operations; clauses with calls, conditionals, or
nested reductions may fall back to scalar or hybrid execution. Debug and
profiling flags in the CLI expose vectorization counts and reduction-path
decisions.

\begin{table}[htbp]
\centering
\caption{Backend execution strategies for lowered tensor clauses.}
\label{tab:tensor-execution-strategies}
\renewcommand{\arraystretch}{1.08}
\begin{tabular}{@{}L{0.24\textwidth}L{0.64\textwidth}@{}}
\toprule
Strategy & Used when \\
\midrule
Vectorized & Clause body can be evaluated with NumPy broadcasting over all loop axes. \\
Scalar & Clause body depends on features that are safest to evaluate one output point at a time. \\
Hybrid & Some axes are scalarized while remaining axes are evaluated with broadcast arrays. \\
Call-scalar & Function calls depend on loop variables in a way that prevents full vectorization. \\
Specialized reduction & A lowered reduction matches a known kernel pattern such as matrix-like contraction. \\
\bottomrule
\end{tabular}
\end{table}

\section{Recurrence Dependency Extraction}

Recurrence analysis is shared between compiler and backend helpers in
\pathcode{src/einlang/shared/recurrence_analysis.py}. The analysis searches
lowered clause bodies for reads of the same binding DefId that is being written.
It then checks whether each read subscript is the current loop variable, a
constant offset from the loop variable, or a reduction variable bounded by the
loop variable.

For a recurrent binding over indices $(i_1,\ldots,i_n)$, a self-read such as
\texttt{H[i - 1, j]} contributes offset $(-1,0)$ and a read such as
\texttt{H[i, j - 1]} contributes $(0,-1)$. These offsets define a point
dependency graph. A legal execution order is any topological order of that
graph. For common non-positive offsets, row-major or timestep-major traversal is
legal after boundary clauses have been checked.

\subsection{Recurrence Detection Algorithm}

Recurrence detection is structural. A clause is recurrent when its body reads
the same binding DefId that the surrounding lowered Einstein writes, and at
least one read index is not the same output loop variable at that dimension.

\begin{lstlisting}[style=pythoncode,caption={Structural recurrence detection.}]
recurrence_dims(clause, written_defid):
    loop_dims = map clause output dimensions to clause loop positions
    loop_defids = [loop.variable.defid for loop in clause.loops]
    read_index_lists = collect_rectangular_reads(clause.body, written_defid)
    out = []

    for loop_pos, output_dim in enumerate(loop_dims):
        loop_defid = loop_defids[loop_pos]
        for read_indices in read_index_lists:
            if output_dim >= len(read_indices):
                continue
            read_index = read_indices[output_dim]
            if read_index is exactly loop_defid:
                continue
            if read_index is loop_defid - positive_literal:
                out.append(loop_pos)
                break
            if read_index is bounded reduction variable tied to loop_defid:
                out.append(loop_pos)
                break
            if read_index is same timestep but another clause writes first:
                out.append(loop_pos)
                break
    return out
\end{lstlisting}

The same pass also handles same-timestep dependencies. If one clause writes
\texttt{state[t, 0]} and a later clause reads \texttt{state[t, 0]} while writing
\texttt{state[t, 1]}, the second clause must run inside a timestep-major loop
even though the read is not \texttt{t - 1}. This is why the compiler annotates
all lowered clauses with \texttt{recurrence\_dims\_override}; the backend can
then use the planned order directly.

The recurrence detector is structural because the source language exposes
self-reads directly. It does not need to infer a dependency graph from arbitrary
Python control flow. Instead, it searches for reads of the same written DefId
and compares the read index expression with the current loop variable. Exact
use means same point, negative literal offsets mean previous points, and
bounded reduction variables can represent dependencies over a finite prefix.
Those cases cover the current optimized recurrence fragment.

The conservative cases are just as important. If a read expression is dynamic
or cannot be related to the loop variable, the compiler should not invent a
finite ordering or storage window. It can still execute by materializing the
full tensor, but it should not claim an optimization proof. This is why the
analysis returns dimensions and overrides rather than rewriting the program
into a fixed loop schedule too early.

The detector's unit of identity is the written binding DefId. This prevents a
read of another value with the same printed name in a different scope from being
mistaken for a self-dependency. It also lets recurrent dependencies survive
function and module elaboration: the recurrence question is "does this clause
read the binding it is writing?" rather than "does this text contain the same
identifier spelling?" That is the same identity discipline used by range and
shape analysis, now applied to temporal dependencies.

Same-timestep dependencies are the non-obvious case. A clause can depend on a
value written earlier at the same timestep rather than on a previous timestep.
This is not a negative offset, but it still constrains order. Treating that case
as recurrence metadata rather than as a separate backend trick lets the compiler
choose timestep-major execution when needed. The algorithm therefore captures
both historical dependencies and intra-step dependencies in one annotation
scheme, which is necessary for dynamic-programming kernels that update several
fields per timestep.

\textit{Concrete example.} For \texttt{let fib[n in 2..N] = fib[n - 1] +
fib[n - 2];}, the detector finds two reads of the written DefId and extracts
backward offsets one and two along the only output dimension. The recurrence
dimension is therefore dimension zero, and storage inference later knows the
lookback is two. For a two-field state where \texttt{state[t,1]} reads
\texttt{state[t,0]}, the read has no negative offset, but it still constrains
clause order inside each \texttt{t}; the detector records a same-timestep
dependency so the backend executes the producer clause first.

\section{Recurrence Ordering}

\pathcode{src/einlang/passes/recurrence_order.py} annotates lowered Einstein
clauses with recurrence dimensions. It handles ordinary backward dependencies,
hybrid dependencies where reduction variables are bounded by loop variables, and
same-timestep dependencies where one clause writes a point that another clause
reads at the same recurrence index. It can override recurrence dimension order
for cases such as Cholesky-like programs where the natural dependency order is
not simply the textual loop order.

\section{Storage Window Inference}

The same offset information drives storage decisions. If a one-dimensional
recurrence only reads the previous two values and later code reads the final
value, the runtime does not need the full sequence. It needs enough history for
self-reads and enough suffix for later observations:

\[
\mathit{preserve} = \max(\mathit{lookback}, \mathit{tail})
\]

In the implementation, this calculation is attached to the lowered recurrence
node as storage metadata: whether the output must be fully materialized, which
output dimension is the recurrence dimension, how many steps must be preserved,
and whether downstream observations are bounded. The runtime consults this
metadata when deciding whether to allocate a full tensor or a circular buffer.

Whole-tensor observations and dynamic tail reads force full materialization. The
current implementation demonstrates finite-window execution for the primary
recurrence axis, including multidimensional tensor values whose non-windowed
axes are retained fully. The circular buffer implementation in
\pathcode{src/einlang/backends/numpy_einstein_mixin_setup.py} can also represent
multiple circular dimensions, which prepares the runtime for general
multidimensional vector windows.

\subsection{Circular Buffer Execution}

The circular recurrence buffer stores only the preserved suffix along the
recurrence axis. Writes are remapped by modulo arithmetic; reads of retained
history are remapped to the corresponding buffer slot. If a value later escapes
as a public output and the compiler did not mark it as safely lazy, the backend
materializes the buffer into a full tensor before returning it. This keeps the
observable output model simple while allowing local recurrence computations to
use less memory.

\begin{lstlisting}[style=pythoncode,caption={Finite-window recurrence-storage inference.}]
infer_storage(lowered_recurrence, later_statements):
    recurrence_dim = first recurrence output dimension
    recurrence_extent = static_extent(lowered_recurrence.shape[recurrence_dim])

    lookback = 0
    for recurrent_clause in lowered_recurrence.body.items:
        for self_read in reads_of_written_defid(recurrent_clause.body):
            offset = backward_offset(self_read.indices[recurrence_dim])
            if offset is unknown:
                return requires_full_output
            lookback = max(lookback, offset)

    tail = 0
    for later in later_statements:
        for use in uses_of_written_defid(later):
            if use is whole tensor or dynamic index:
                return requires_full_output
            idx = static_int(use.indices[recurrence_dim])
            tail = max(tail, recurrence_extent - idx)

    history_window = lookback + 1
    preserve_steps = max(history_window, tail)
    return circular_buffer(axis=recurrence_dim, preserve_steps=preserve_steps)

execute_circular_recurrence(node):
    if node.requires_full_output or preserve_steps >= full_extent:
        execute full-history recurrence
    ring = CircularRecurrenceBuffer(shape, axis, preserve_steps)
    bind output DefId to ring in the environment
    execute initial clauses into ring
    for timestep in recurrence_loop:
        execute recurrent clauses for this timestep
        map every write/read on recurrence axis through index % preserve_steps
        if unsupported observation materializes ring:
            fall back to full materialization
    return ring
\end{lstlisting}

The storage algorithm combines two independent obligations. The first is
history: recurrent clauses need enough previous timesteps to satisfy self-reads
such as \texttt{t - 1} or \texttt{t - 2}. The second is observation: later code
may request the final value, the last few values, or the whole tensor. The
preserved window is the maximum of those requirements. If either side cannot be
bounded statically, the compiler chooses full materialization.

The execution half is written to preserve observable semantics. The environment
binds the output DefId to a circular buffer so recurrent self-reads and writes
use the same identity as a full tensor would. Internally, the recurrence axis is
remapped through modulo arithmetic. Externally, if a value must escape as a
normal array and the compiler has not proven that a suffix is enough, the ring
is materialized before returning. The optimization is therefore invisible to
the source program except through memory use and profiling.

The current implementation focuses on a primary recurrence axis, but the data
structure anticipates more general windows. This is why the algorithm keeps the
axis and preserved step count explicit in metadata instead of baking a
one-dimensional special case into every backend path. Future multidimensional
storage inference can extend the proof without replacing the execution model.

The preservation formula deliberately includes the current timestep by using
\texttt{lookback + 1}. A recurrence that reads \texttt{t - 2} needs the values
for the current write plus two earlier slots to make modulo indexing
unambiguous. The tail term is separate because later observations impose a
different constraint: a program that reads the final three values needs those
three values even if the recurrence itself only reads one step back. Taking the
maximum of the two obligations is what lets the compiler use one buffer for both
internal recurrence correctness and external observation.

The fallback cases protect the source-level value model. Whole-tensor reads,
dynamic indices, and unknown backward offsets all force full materialization
because the compiler cannot prove which points will be observed. This is an
intentional asymmetry: the implementation can always allocate more memory and
preserve semantics, but it must not allocate less memory unless the proof is
available. The thesis therefore treats circular buffering as a certified
optimization. The metadata attached to the lowered recurrence is the
certificate, and the runtime optimization is allowed only when the certificate
states a bounded window.

\textit{Concrete example.} For the Fibonacci program above with a later read
\texttt{fib[N - 1]}, the recurrent clause needs two previous values, so
\texttt{history\_window = 3}; the observation needs only the final value, so
\texttt{tail = 1}; the chosen circular buffer preserves three steps. For a
program that later returns \texttt{fib} as a whole tensor, the tail cannot be
bounded to a suffix, so the same recurrence is fully materialized. The numerical
recurrence is identical in both cases; only the proven observation pattern
changes the storage plan.

\begin{table}[htbp]
\centering
\caption{Representative recurrence-storage cases covered by tests.}
\label{tab:storage-checks}
\renewcommand{\arraystretch}{1.06}
\begin{tabular}{@{}lcccl@{}}
\toprule
Case & lookback & tail & preserve & Decision \\
\midrule
two-step, final & 2 & 1 & 3 & circular buffer \\
one-step, last 3 & 1 & 3 & 3 & circular buffer \\
symbolic T, final & 1 & 1 & 2 & circular buffer \\
multidimensional tensor & 1 & 1 & 2 & circular buffer on recurrence axis \\
whole tensor & 1 & -- & -- & full materialization \\
dynamic tail & 1 & -- & -- & full materialization \\
\bottomrule
\end{tabular}
\end{table}

\section{Current Boundary}

The current compiler proves finite storage only for static-offset recurrence
patterns with an inferred primary recurrence axis. It falls back to full
materialization when source observations demand the whole tensor or when the
tail bound cannot be statically proven. This conservative boundary is important:
the implementation should miss some optimization opportunities before it risks
returning an incomplete recurrence value.
